\subsection{Common core of the \(3\) formulations of T-duality}

Let us define the stable domain
\[
 \mathscr D_0=\mathbb Z^2\times\mathbb R_{>0},
 \qquad
 E_{R_0}(m,w)=\frac{m^2}{R_0^2}+w^2R_0^2,
\]
and the involution
\[
 \mathsf T_0(m,w,R_0)=(w,m,R_0^{-1}).
\]

\begin{theorem}[Residual ambient core equivalence--quotient]
\label{thm:cierre-vtm-tres-formulaciones-t}
After prolonging the variables \(r,Q\) to the stable domain
\(\mathbb Z^2\), the three formulations are identified by means of
\[
 I_F(r,Q,R)=(m,w,R_0)=(r,Q,R),
 \qquad I_{16\to02}=\operatorname{id}_{\mathscr D_0}.
\]
These maps intertwine the involutions and the quadratic forms:
\[
 I_F\mathsf T_F=\mathsf T_0I_F,
 \qquad
 E_R^F=E_{R_0}\circ I_F,
 \qquad
 \mathsf T_0^2=\operatorname{id},
 \qquad
 E_{R_0}=E_{R_0^{-1}}\circ\operatorname{swap}.
\]
\end{theorem}

\begin{proof}
The two maps are renamings of the same three coordinates. Moreover,
\[
 \mathsf T_0^2(m,w,R_0)=(m,w,R_0)
\]
and
\[
 E_{R_0^{-1}}(w,m)
 =\frac{w^2}{R_0^{-2}}+m^2R_0^{-2}
 =w^2R_0^2+\frac{m^2}{R_0^2}
 =E_{R_0}(m,w).
\]
The identities are the same after applying \(I_F\) or the identity between the
two formulations.
\end{proof}

In \(\mathcal H_{\mathrm{lat}}=\ell^2(\mathbb Z^2)\), let us define
\[
 H(R_0)|m,w\rangle=E_{R_0}(m,w)|m,w\rangle,
 \qquad U_T|m,w\rangle=|w,m\rangle.
\]
Then \(U_T\) is unitary, \(U_T^2=I\), and the preceding scalar identity
is equivalent, on the orthonormal basis, to
\[
 \boxed{U_TH(R_0)U_T^{-1}=H(R_0^{-1}).}
\]
This is the second exact realization of the same core, not a new physical hypothesis.

\subsection{Distinction of the 2 residual sections by their energy insertions and equivalence classes}

The positive residual convention uses
\[
 r_9^+(n)=1+((n-1)\bmod9),
 \qquad Q_9^+(n)=\frac{n-r_9^+(n)}9,
\]
whereas the zero residual convention inserts the class
\(\bar r_9(n)\in\{0,\ldots,8\}\). The change of section that preserves the
integer is
\[
 C(r,Q)=\bigl(\bar r,Q+\mathbf1_{\{r=9\}}\bigr).
\]

\begin{proposition}[Energy Obstruction of Section Change]
\label{prop:cierre-vtm-obstruccion-secciones-nonadicas}
The map \(C\) preserves \(n=9Q+r\), but in general it preserves
neither the form \(E_R\) nor the exchange involution. Equivalence of
the ambient kernel therefore does not imply equivalence of the two nonadic
embeddings.
\end{proposition}

\begin{proof}
For \(n=9\), the first convention gives \((r,Q)=(9,0)\), whereas the change with carry gives
\((\bar r,Q_0)=(0,1)\). Their energies are
\[
 E_R(9,0)=\frac{81}{R^2},
 \qquad E_R(0,1)=R^2,
\]
which do not coincide for every \(R>0\). Insertion without correcting
the carry would send this example to \((0,0)\) and would likewise fail to
preserve the integer. A unique section must be fixed, or a form covariant
under change of section must be constructed, before this equivalence is
promoted.
\end{proof}

\begin{definition}[Quotient of section change and covariant energy]
\label{def:cierre-vtm-energia-covariante-seccion}
Let
\[
 L_9:=\mathbb Z(9,-1)\subset\mathbb Z^2,
 \qquad
 S(x_1,x_2)=(x_2,x_1),
\]
and let
\[
 q_R(x_1,x_2)=\frac{x_1^2}{R^2}+x_2^2R^2.
\]
On the quotients
\[
 \mathscr Q_9=(\mathbb Z^2/L_9)\times\mathbb R_{>0},
 \qquad
 \mathscr Q_9^\vee=(\mathbb Z^2/SL_9)\times\mathbb R_{>0},
\]
we define
\[
 \mathcal E_R^{L_9}([x]_{L_9})
 :=\min_{\ell\in L_9}q_R(x+\ell)
\]
and
\[
 \mathsf T_{\mathrm{cov}}([x]_{L_9},R)
 :=([Sx]_{SL_9},R^{-1}).
\]
\end{definition}

\begin{theorem}[Covariant Equivalence of Nonary Sections]
\label{thm:cierre-vtm-equivalencia-covariante-secciones}
The energy \(\mathcal E_R^{L_9}\) and the transformation
\(\mathsf T_{\mathrm{cov}}\) are well defined. The conventions
\((9,0)\) and \((0,1)\), and in general the sections \(1,\ldots,9\) and
\(0,\ldots,8\) with the corresponding carry, represent the same class
of \(\mathbb Z^2/L_9\). Moreover,
\[
 \boxed{
 \mathcal E_{R^{-1}}^{SL_9}([Sx]_{SL_9})
 =\mathcal E_R^{L_9}([x]_{L_9}).}
\]
The inverse exchange, defined on \(\mathscr Q_9^\vee\), returns the
initial class; hence the two quotiented charts are T-dual.
\end{theorem}

\begin{proof}
If \(x'=x+\ell_0\), the translation
\(\ell\mapsto\ell+\ell_0\) permutes \(L_9\), so the minimum does not depend on the
representative. It exists because \(q_R(x+k(9,-1))\) is a coercive quadratic function
of \(k\in\mathbb Z\). The difference
\((9,0)-(0,1)=(9,-1)\) belongs to \(L_9\), and the same occurs under every change
of section with carry. Finally,
\[
\begin{aligned}
 \mathcal E_{R^{-1}}^{SL_9}([Sx])
 &=\min_{\ell\in L_9}q_{R^{-1}}(Sx+S\ell)\\
 &=\min_{\ell\in L_9}q_R(x+\ell)
 =\mathcal E_R^{L_9}([x]).
\end{aligned}
\]
Since \(S^2=I\), the second exchange recovers the original quotient and the
radius \(R\).
\end{proof}

\paragraph{Equivalence domain of the dual nucleus and nonequivalence of energy insertions.}
The equivalence of the kernels over \(\mathscr D_0\) is exact. The two
embeddings are not equivalent for the uncorrected energy; the covariant
closure is the quotient and the energy of the preceding theorem.

\subsection{Exact normalization of the string realization}

Fix \(\alpha'>0\). In the domain
\(\mathscr D_{\alpha'}=\mathbb Z^2\times\mathbb R_{>0}\), let
\[
 \mathsf T_{\alpha'}(m,w,R)=\left(w,m,\frac{\alpha'}R\right),
 \qquad
 \mathcal N_{\alpha'}(m,w,R)
 =\left(m,w,\frac R{\sqrt{\alpha'}}\right).
\]

\begin{theorem}[String scale conjugation]
\label{thm:cierre-vtm-normalizacion-alpha-t}
The normalization \(\mathcal N_{\alpha'}\) is an exact conjugation:
\[
 \boxed{\mathcal N_{\alpha'}\mathsf T_{\alpha'}
 =\mathsf T_0\mathcal N_{\alpha'}.}
\]
For
\[
 E_{\mathrm{str}}(m,w,R)
 =\frac{m^2}{R^2}+\frac{w^2R^2}{\alpha'^2},
\]
we have
\[
 E_{\mathrm{str}}
 =\alpha'^{-1}E_0\circ\mathcal N_{\alpha'}.
\]
Furthermore, if
\[
 p_L=\frac mR+\frac{wR}{\alpha'},
 \qquad p_R=\frac mR-\frac{wR}{\alpha'},
\]
then \(p_L\) remains fixed and \(p_R\mapsto-p_R\).
\end{theorem}

\begin{proof}
We have
\[
 \mathcal N_{\alpha'}\mathsf T_{\alpha'}(m,w,R)
 =\left(w,m,\frac{\sqrt{\alpha'}}R\right)
 =\mathsf T_0\mathcal N_{\alpha'}(m,w,R).
\]
The energy identity results from
\[
 E_0\!\left(m,w,\frac R{\sqrt{\alpha'}}\right)
 =\frac{\alpha'm^2}{R^2}+\frac{w^2R^2}{\alpha'}
 =\alpha' E_{\mathrm{str}}(m,w,R).
\]
The substitution of
\((m,w,R)\mapsto(w,m,\alpha'/R)\) in \((p_L,p_R)\) gives, respectively,
\(p_L\) and \(-p_R\).
\end{proof}

Conjugation and spectral invariance with fixed oscillatory data are exact.
The identification of \(m,w,R,\alpha'\) with momentum, winding, physical radius,
and inverse tension is a physical interface with its own premises.

\subsection{Level, linear mass exchange, and modularity: \(3\) nonfusion controls}

The active string condition is
\[
 N_L-N_R=a_L-a_R-mw.
\]
The schematic formula \(N_L-N_R=nw\) coincides with it only
after declaring \(a_L=a_R\) and \(n=-m\), or an equivalent reversal of
leaf orientation. Without that dictionary, the two expressions are not the
same formulation.

The linear symmetry of the action is typified in the extended space of
coefficients and modes:
\[
 L(A,C;n,w)=nA+wC,
 \qquad
 \sigma(A,C;n,w)=(C,A;w,n).
\]
Then
\[
 L\circ\sigma=wC+nA=L.
\]
It is an exact identity under simultaneous exchange of parameters and
coordinates. It is not equivalent to quadratic invariance of \(H(R)\) on a
sector with fixed physical parameters; such a promotion requires an
additional intertwiner.

Finally, on the imaginary axis of the upper half-plane, let us define
\[
 \iota(R_0)=iR_0^2,
 \qquad S(\tau)=-\frac1\tau.
\]
Then
\[
 \boxed{S\circ\iota(R_0)=\iota(R_0^{-1}).}
\]
Since \(J(S\tau)=J(\tau)\), one obtains
\(J(iR_0^2)=J(iR_0^{-2})\). This is an exact semiconjugacy of the radius
parameter with the modular transformation \(S\), proved from the explicit
starting modular structure. It does not by itself identify target-space
duality, the exchange of charges, or the construction of
\(V^\natural\).
